\chapter{共线性诊断与岭回归}

\begin{studysource}
\textbf{来源说明}　课件：第二讲第 10–15 页　课堂代码：回归诊断.R 末尾　脚本：\texttt{scripts/05\_共线性与岭回归.R}　数据：datasetB
\end{studysource}


\section{后果：五个模型并排}
\begin{rcode}[R · 同一批人、不同自变量组合]
x <- c("age", "bmi", "waist")
round(cor(d[, c(x, "glu")]), 3)                  # r(bmi, waist) = 0.968
mods <- list(bmi = lm(glu ~ bmi, d), waist = lm(glu ~ waist, d),
             "age+bmi" = lm(glu ~ age + bmi, d), "age+waist" = lm(glu ~ age + waist, d),
             "age+bmi+waist" = m)
get <- function(fit, v) {                        # 取某个变量的 b、SE、P，不在模型里就给 NA
  out <- if (v %in% names(coef(fit))) coef(summary(fit))[v, c(1, 2, 4)] else rep(NA, 3)
  setNames(out, c("b", "SE", "P"))
}
round(t(sapply(mods, function(f) c(bmi = get(f, "bmi"), waist = get(f, "waist")))), 4)
\end{rcode}
\begin{routput}
> round(t(sapply(mods, ...)), 4)
               bmi.b bmi.SE  bmi.P waist.b waist.SE waist.P
bmi           0.0634 0.0179 0.0005      NA       NA      NA
waist             NA     NA     NA  0.0266   0.0072  0.0003
age+bmi       0.0520 0.0138 0.0002      NA       NA      NA
age+waist         NA     NA     NA  0.0200   0.0056  0.0005
age+bmi+waist 0.0671 0.0552 0.2260 -0.0063   0.0223  0.7780
\end{routput}
单独放都显著，一起放都不显著，waist 还变了号；但整个方程的 $R^2$ 和 $F$ 几乎不受影响。共线性伤的是“解释单个系数”，不是预测。

\section{VIF、容忍度与条件数}
\begin{rcode}[R · 共线性诊断]
vif(m)                                           # 1.03、16.02、16.14
1 / vif(m)                                       # 容忍度
1 / (1 - summary(lm(bmi ~ age + waist, data = d))$r.squared)   # 手算 bmi 的 VIF
sapply(x, function(v)                            # 一行算全部
  1 / (1 - summary(lm(reformulate(setdiff(x, v), v), data = d))$r.squared))
coef(summary(m))["bmi", 2] / coef(summary(lm(glu ~ age + bmi, d)))["bmi", 2]   # 4.00
sqrt(vif(m)["bmi"])                              # 4.00：√VIF 就是标准误膨胀倍数
ev <- eigen(cor(d[, x]))$values                  # 1.988、0.981、0.032
sqrt(max(ev) / min(ev))                          # 条件数 kappa = 7.93
kappa(scale(d[, x]), exact = TRUE)               # 同一个数
\end{rcode}
VIF 的经验界是 5（警惕）和 10（严重），条件数的界在教材里从 15 到 30 不等。本例 bmi 加进 waist 后标准误恰好放大 4.00 倍，等于 $\sqrt{16.02}$。


\begin{studyexplain}
\textbf{课堂代码里的两个小问题}


\Verb|回归诊断.R| 末尾写的是 \Verb|vif(fit)|，而此时 \Verb|fit| 已被前一行改成了 \Verb|regsubsets| 对象，会报错 \Verb|缺少参数"id"|；\Verb|vif()| 只接受 lm 对象。另一行 \Verb|kappa(scale(d))| 把因变量 glu 也放进了矩阵，而且 \Verb|kappa()| 默认只给近似值。诊断共线性只放自变量，并写 \Verb|exact = TRUE|。

\end{studyexplain}


\section{岭回归：MASS::lm.ridge()}
\begin{rcode}[R · 拟合一串 lambda]
lam <- seq(0, 100, by = 0.5)
rr  <- lm.ridge(glu ~ age + bmi + waist, data = d, lambda = lam)
select(rr)                                  # HKB 1.16、L-W 1.21、GCV 最小在 7.5
rr$lambda[which.min(rr$GCV)]                # 7.5
round(t(rr$coef)[lam %in% c(0, 7.5, 20, 30, 60), ], 3)    # 画岭迹用的系数
round(coef(rr)[lam %in% c(0, 7.5, 30), ], 4)              # 写方程用的系数
\end{rcode}
\begin{routput}
> round(t(rr$coef)[lam %in% c(0, 7.5, 20, 30, 60), ], 3)
        age   bmi  waist
  0.0 0.454 0.216 -0.050
  7.5 0.429 0.129  0.035
 20.0 0.397 0.100  0.061
 30.0 0.375 0.091  0.066
 60.0 0.322 0.079  0.070
> round(coef(rr)[lam %in% c(0, 7.5, 30), ], 4)
                age    bmi   waist
  0.0 2.4617 0.0314 0.0671 -0.0063
  7.5 2.3235 0.0297 0.0400  0.0044
 30.0 2.4546 0.0260 0.0283  0.0083
\end{routput}


\begin{studyexplain}
\textbf{lm.ridge 的两套系数}


\Verb|rr$coef| 是“变量 × λ”的矩阵，自变量已经标准化（按 $n$ 作分母的 SD），不含截距，画岭迹用它，画之前要 \Verb|t()| 转置；\Verb|coef(rr)| 是“λ × 变量”的矩阵，已还原到原始单位并带截距，写回归方程用它。$\lambda=0$ 那一行就是 OLS。

\end{studyexplain}


\begin{rcode}[R · 岭迹图]
matplot(rr$lambda, t(rr$coef), type = "l", lty = 1, lwd = 2,
        col = c("#0072B2", "#D55E00", "#009E73"),
        xlab = expression(lambda), ylab = "standardized coefficient")
abline(h = 0, col = "grey70"); abline(v = 7.5, lty = 2)
legend("topright", x, col = c("#0072B2", "#D55E00", "#009E73"), lty = 1, bty = "n")
plot(rr$lambda, rr$GCV, type = "l")         # GCV 曲线
\end{rcode}


\includegraphics{fig2_ridge.pdf}

\figtitle{图 2　datasetB 的岭迹与 GCV 曲线。waist 系数在 λ 很小时就穿过 0 变正，之后稳定；GCV 在 0–20 之间几乎是平的}


GCV 在 $\lambda=0$ 与 $\lambda=7.5$ 之间只差 0.6\%（图 \ref{fig:ridge}b）。GCV 衡量的是预测误差，而共线性不伤预测，所以它帮不上“读系数”的忙；岭参数要结合岭迹是否平稳、系数是否合理来定，并说明理由。

\section{自己写一遍岭回归}
\begin{rcode}[R · 按定义手算岭回归]
lambda <- 30
X  <- as.matrix(d[, x])
Xc <- scale(X, center = TRUE, scale = FALSE)
sc <- sqrt(colSums(Xc^2) / nrow(X))        # lm.ridge 用 n 作分母
Z  <- sweep(Xc, 2, sc, "/")                # 标准化后的自变量
b_std <- solve(crossprod(Z) + lambda * diag(3), crossprod(Z, d$glu - mean(d$glu)))
b_raw <- drop(b_std) / sc                  # 还原到原始单位
b0    <- mean(d$glu) - sum(b_raw * colMeans(X))
c(b0, b_raw)                               # 与 coef(lm.ridge(..., lambda = 30)) 完全相同
R <- cor(X); ry <- cor(X, d$glu)
solve(R + (lambda / nrow(X)) * diag(3), ry)   # 课件笔记的 (R + kI)^{-1} r_y，k = lambda/n
\end{rcode}
最后一行与课件笔记的相关阵形式一致，只是尺度不同：\Verb|lm.ridge| 的 $\lambda$ 对应 $k=\lambda/n$。不同软件对标准化的分母（$n$ 或 $n-1$）和 $\lambda$ 的归一方式不同，岭参数的数值不能跨软件直接比。

\section{比岭回归更常用的处理}
\begin{rcode}[R · 只留一个，或合成一个]
summary(lm(glu ~ age + bmi, data = d))                       # 只留 BMI
d$obesity <- (as.numeric(scale(d$bmi)) + as.numeric(scale(d$waist))) / 2
summary(lm(glu ~ age + obesity, data = d))                   # z 分数平均
pc <- prcomp(d[, c("bmi", "waist")], scale. = TRUE)
summary(pc)                                                  # 第一主成分占 98.4%
\end{rcode}
课件第 15 页给的顺序是：先由研究问题决定只留一个，其次合成一个指标，再次才是岭回归。合成指标与第一主成分在本例给出相同的 $t$（3.70），因为两个标准化变量的平均就是第一主成分的方向（主成分的正负号是任意的）。

\section{强影响点会改变岭参数的选择}
\begin{routput}
> for (s in c("all", "del149")) ...
all : GCV 最优 lambda = 7.5 ；bmi、waist 同时为正的最小 lambda = 3
del149 : GCV 最优 lambda = 1.5 ；bmi、waist 同时为正的最小 lambda = 18.5
\end{routput}
删掉 149 号后，要把 waist 拉成正号需要大得多的 λ。作业二第二题因此在“全样本、删 149、删 150、两人都删”四种样本下画岭迹，选了一个在四种情形下都已平稳的 λ = 30，而没有采用全样本 GCV 给出的 7.5（第 9 章）。
